\begin{table*}[t]
\centering
% \small
\setlength{\tabcolsep}{6pt}
\renewcommand{\arraystretch}{1.0}
\caption{Main results. For latency, we show the total inference time for $128$ issues.}
\resizebox{0.9\textwidth}{!}{
\begin{tabular}{c c c c cc cc}
\toprule
\multirow{2}{*}{Agent} & \multirow{2}{*}{Memory} & \multirow{2}{*}{\%Resolved} &
\multirow{2}{*}{\#Tool Calls} &
\multicolumn{2}{c}{w/o CPU Offload} &
\multicolumn{2}{c}{w/ CPU Offload} \\
& & &  &
Latency ($\times 128$/s) & Speedup &
Latency ($\times 128$/s) & Speedup \\
\midrule
\multirow{4}{*}{\textbf{\shortstack{DeepSWE\\(32B Dense)}}}
 & Full-Context   & 40.4 &  34.39 & 9457.78 & 1$\times$ & 6390.62 & 1$\times$ \\
 \arrayrulecolor{lightgray}\cmidrule[0.5pt]{2-8}\arrayrulecolor{black}
 & Qwen3-4B (base) & 29.9 & 24.74 & 4360.72 & 2.17$\times$ & 3649.61 & 1.75$\times$ \\
 & Qwen3-4B (SFT) & 39.8 & 35.61 & 5232.64 & 1.81$\times$ & 4168.85 & 1.53$\times$\\
 & $\textbf{GRPO}_{AC}$ & \textbf{41.0} & 40.69 & 5622.33 & 1.68$\times$ & 4400.89 & 1.45$\times$ \\
\midrule
\multirow{5}{*}{\textbf{\shortstack{Qwen3-Coder-Max\\(480B A35B)}}}
 & Full-Context   & \textbf{57.2} & 36.68 & 6291.11 & 1$\times$ & 5129.90 & 1$\times$ \\
 \arrayrulecolor{lightgray}\cmidrule[0.5pt]{2-8}\arrayrulecolor{black}
 & No Summary  & 46.6 & 25.20 & 3780.74 & 1.66$\times$ & 3421.78 & 1.50$\times$\\
 & Qwen3-4B (base) & 42.2 & 39.40 & 3421.63 & 1.84$\times$ & 3138.97 & 1.63$\times$\\
 & Qwen3-4B (SFT) & 47.3 & 43.39 & 3819.70 & 1.65$\times$ & 3337.58 & 1.54$\times$ \\
 & $\textbf{GRPO}_{AC}$ & 51.0 & 44.55 & 3917.68 & 1.61$\times$ & 3594.51 & 1.43$\times$ \\
\midrule
\multirow{5}{*}{\textbf{\shortstack{GLM-4.7\\(355B 32B)}}}
 & Full-Context   & \textbf{69.0} &  55.11 & 6361.49 & 1$\times$ & 5869.42 & 1$\times$ \\
 \arrayrulecolor{lightgray}\cmidrule[0.5pt]{2-8}\arrayrulecolor{black}
 & No Summary  & 59.4 & 46.82 & 4216.83 & 1.51$\times$ & 3842.17 & 1.53$\times$ \\
 & Qwen3-4B (base) & 58.3 & 49.32 & 3928.44 & 1.62$\times$ & 3526.91 & 1.66$\times$\\
 & Qwen3-4B (SFT) & 61.3 & 52.74 & 4087.92 & 1.56$\times$ & 3668.53 & 1.60$\times$ \\
 & $\textbf{GRPO}_{AC}$ & 62.7 & 51.21 & 3318.42 & $1.92\times$ & 2821.38 & 2.08$\times$ \\
\bottomrule
\end{tabular}
}
\label{tab:main}
\vspace{-3mm}
\end{table*}

% The summarization capability of base model might be limited based on our empirical experimentation. Thus we propose the following training pipeline to align the performance of summarization to the full-context baseline.
% Notice that, our training pipeline assumes the agent model to be freezed, and thus efficient without any online environment interactions.
% Specifically, we first generate trajectories from the standard agent pipeline. Then we collect data by treating the interaction histories at each step as input and the actual generated action as the desired output. We then perform GRPO training on these data.
% The novel part is the reward design, where we match the output action from agent given summaries and recent turns as with the expected action. We use the similarity between these two actions as a reward signal for our GRPO training. And thus the model can capture the sufficient statistics so from the history.
Empirical analysis suggests that the inherent summarization capability of base models is often insufficient to compress the history for complex agentic interactions, especially with a dedicated compression ratio.
To bridge this gap, we introduce a training pipeline designed to align the summarization model's output with what the agent requires for optimal decision-making.
The key idea is to train the memory model to capture the \emph{sufficient statistics} of the interaction history---the minimal representation from which the agent can recover behavior equivalent to full-context inference~\cite{sutton1998reinforcement}. Formally, a summary $s = f(\tau_{<t-k})$ is sufficient if the agent's action conditioned on the summary matches what it would produce given the full history: $\pi(\cdot \mid s, C_t) \approx \pi(\cdot \mid \tau_{<t})$, where $C_t = \tau_{t-k:t-1}$ denotes the most recent $k$ raw turns. Rather than optimizing for generic text quality, we directly optimize for this \emph{functional equivalence}---whether the compressed memory induces correct downstream behavior.

Crucially, the agent model $\pi$ remains frozen throughout training. This decouples the optimization of the memory model $f$ from the agent's reasoning policy, enabling an efficient offline regime that requires no online environment interactions. Concretely, we first generate reference trajectories using the standard full-context agent pipeline. For each step $t$ in a trajectory, we record the ground-truth action $a_t^* = \pi(\cdot \mid \tau_{<t})$ produced by the full-context agent.
We then fine-tune the memory model using GRPO (Eq.~\ref{eq:grpo_loss}) with the following reward: given a candidate summary $s = f(\tau_{<t-k})$, the frozen agent produces $\hat{a}_t = \pi(\cdot \mid s, C_t)$, and we define
\begin{equation}
R(s) = \text{sim}(\hat{a}_t,\; a_t^*)
\end{equation}
By maximizing this reward while enforcing a clipping on the maximum response length, the memory model learns to distill precisely the information the agent needs for decision-making, discarding irrelevant details that would otherwise consume context budget.